% Adjoining prose replacements, shown outside their chapter positions.
The full bounded maximum principle requires capacity and is proved in
\Cref{thm:bounded-maximum-principle}.  Here we retain real model inputs,
but must distinguish them from rational model inputs: the maximum of
real model functions can be continuous without being a model function.
The measure of that maximum is already available from the bounded
construction and \Cref{prop:ch24-continuous-uniform}.

On the open set where \(\varphi\) strictly exceeds \(\psi\), the measure
of the maximum should see only \(\varphi\).  For rational potentials we
prove this by choosing a common SNC determination; the resulting
measures are finite sums of Dirac masses, and their coefficients are
intersection numbers, which can be real because \(\theta\) is real.
The geometric mechanism is that the coefficient at \(x_E\) depends only
on the divisor classes restricted to \(E\).  We then pass to real model
inputs by uniform approximation at the fixed form \(\theta\), testing on
compact subsets of the strict inequality locus.  Neither the bounded
locality theorem nor capacity is used in this argument.


\begin{remark}[The mass on the contact set is genuinely new]
\label{rem:ch24-contact-mass}
The strict-locus formula \eqref{eq:model-maximum-principle} does not
determine the measure on the contact set \(\{\varphi=\psi\}\).
Nor does \Cref{cor:ch24-model-locality-open} apply merely because the
potentials agree on a set that is not open.  The measure
\(\MA_\theta(\max\{\varphi,\psi\})\) can have an atom on the contact set
that is present in neither input measure.
\Cref{ex:ch24-locality-check} exhibits this on a chain of three
components, where a new atom of mass \(1\) appears at the crossing point.
The maximum is continuous even when it is not model, and its total
mass is still \(V\) by \Cref{thm:bounded-ma-extension}.
Thus the new atom represents redistribution of mass, not an increase
of total mass.
\end{remark}
